% ============================================================================
% REVISED METHODS (major revision) — paste into the manuscript, mark as changed.
% NeuroFormer architecture is UNCHANGED; revisions are to inputs, training, and
% evaluation. [[...]] = numbers finalized from leak-free runs.
% ============================================================================

\subsection{Overview of Revisions}
The NeuroFormer architecture is unchanged from the original submission: three
modality-specific transformer encoders, a cross-modal fusion transformer with a
\texttt{[CLS]} aggregation token and modality-type embeddings, a classification
head, and optional auxiliary-reconstruction decoders. In response to the reviews,
we revised the \emph{inputs to}, \emph{training of}, and \emph{evaluation of} this
same model: (i) a pretrained EEG foundation-model front-end, (ii) a strict
leak-free subject-level protocol, (iii) modality-dropout training for a
missing-modality regime, (iv) external out-of-distribution validation, and (v) a
calibrated, referral-aware output layer.

\subsection{Foundation-Model EEG Representations}
To address the limited value of the original hand-crafted five-channel EEG
features, we replace them with representations from \textbf{LaBraM}, a pretrained
EEG foundation model~\cite{LaBraM2024}, while keeping the downstream NeuroFormer
EEG encoder identical. Raw MIPDB EEG (128-channel EGI, 500\,Hz) is preprocessed to
the foundation model's specification (0.1--75\,Hz band-pass, 60\,Hz notch,
resampling to 200\,Hz), scaled to microvolts, segmented into non-overlapping
15\,s windows, and clipped at $\pm 800\,\mu$V to bound gross artifacts. A
quality-control gate retains only physiological recordings (median amplitude in
3--150\,$\mu$V, $<$15\% clipped). Each window is mapped to a 200-dimensional frozen
LaBraM embedding via the \texttt{braindecode} \texttt{InterpolatedLaBraM} model,
which projects the electrode montage onto LaBraM's canonical channels. The
resulting per-window embedding sequence is fed to the NeuroFormer EEG encoder
exactly as before (only the input dimensionality changes, $128\!\to\!200$).
Eye-tracking and behavioural features are re-derived from the raw recordings using
the original definitions and validated against the prior pipeline (per-feature
$r\!\approx\!1.0$).

\subsection{Leak-Free Subject-Level Evaluation}
To address concerns regarding fold sensitivity and overparameterization, all
results use \textbf{strict subject-level} partitioning---every recording from a
participant is confined to a single fold---under \textbf{repeated} stratified
$k$-fold cross-validation ($5\times[[R]]$) with bootstrap confidence intervals.
Feature standardization is fit on training folds only. We additionally evaluate a
\textbf{compact} NeuroFormer ($\approx$[[2.9]]M parameters) to match the cohort
size, reporting performance against parameter count.

\subsection{Missing-Modality Training and Asynchronous Fusion}
Because the three modalities are sourced from different paradigms and are not
trial-synchronized, we make the cross-paradigm structure explicit and treat fusion
as \textbf{asynchronous, participant-level} integration rather than synchronized
cross-modal interaction. To reflect real screening conditions---in which modalities
are frequently missing---we train with \textbf{modality dropout}: each modality is
randomly withheld per sample, and the model's modality-type embeddings act as
availability indicators, so a single NeuroFormer serves any subset of
$\{$EEG, eye, behaviour$\}$ with graceful degradation. This design follows
missing-modality clinical-fusion work~\cite{DrFuse2024}.

\subsection{Calibrated, Referral-Aware Prediction}
To support clinical triage, NeuroFormer outputs are calibrated post-hoc with
temperature scaling, and we add split-conformal prediction with risk--coverage
analysis: low-confidence cases abstain and are referred for clinician review. We
report Expected Calibration Error and risk--coverage curves, following conformal
clinical-triage methods~\cite{ConformalTriage2025}.

\subsection{External Validation on HBN}
To assess generalization beyond MIPDB, we externally validate on the independent
\textbf{Healthy Brain Network (HBN)} EEG cohort~\cite{HBNEEG2024}
($N\!=\![[\,]]$, RestingState), harmonized to the same preprocessing. As HBN
provides continuous CBCL psychopathology factors rather than categorical
diagnoses, we evaluate whether the frozen NeuroFormer EEG representation predicts
the general psychopathology factor (regression $r$ and median-split ROC-AUC),
reporting the honest out-of-distribution result and the label-definition
difference.
